\documentclass[fleqn,12pt,a4paper]{article}
\usepackage[utf8]{inputenc}
\usepackage[T1]{fontenc}
\usepackage{lmodern} % Fontes vetoriais
\usepackage[brazilian]{babel}
\usepackage{amsmath,amssymb}
\usepackage{nccmath} % permite centralizar equações mesmo com fleqn
\usepackage{geometry,tabularx,array}
\usepackage{setspace}
\geometry{a4paper,margin=1in}
\onehalfspacing

% TikZ e fluxogramas
\usepackage{booktabs}
\usepackage{caption}
\captionsetup{labelsep=endash,justification=raggedright,singlelinecheck=false}
\usepackage{tikz}
\usetikzlibrary{shapes.geometric, arrows.meta, positioning}

% Ajustes finos de espaçamento
\setlength{\mathindent}{0pt}
\setlength{\abovedisplayskip}{4pt}
\setlength{\belowdisplayskip}{4pt}
\setlength{\abovedisplayshortskip}{0pt}
\setlength{\belowdisplayshortskip}{0pt}
\setlength{\jot}{3pt}

\newcolumntype{L}[1]{>{\raggedright\arraybackslash}p{#1}}
\newcolumntype{R}{>{\raggedright\arraybackslash}X}

% BibTeX estilo ABNT
\usepackage[alf]{abntex2cite}
\renewcommand{\refname}{Referências}


\title{Dimensionamento e Sequenciamento de Lotes de Produção de uma Indústria de Balões de Látex}
\author{}
\date{21/09/2026}

\begin{document}
\maketitle

\newpage

\section{Introdução}

\subsection{Contextualização}
\label{sec:contextualizacao}

O látex de borracha natural chega à indústria de transformação como emulsão aquosa e é convertido em produto por vulcanização, reação irreversível entre o polímero e agentes de cura. A família de artefatos que sai desse processo compartilha elasticidade e impermeabilidade, e se separa pelo regime de especificação a que cada aplicação responde: luvas e preservativos são produtos normalizados, sujeitos a ensaio de barreira lote a lote, enquanto artigos de festa não têm norma de desempenho equivalente e têm valor definido por atributos estéticos. O balão de látex pertence a esse segundo grupo.

Sem especificação funcional que discrimine um balão de outro, a diferenciação é inteiramente de aparência e se decompõe em três atributos. O formato é dado pela forma em que o látex é imerso e fixa contorno e calibre; o acabamento é o tratamento de superfície, fosco, perolado ou brilhante; a cor vem do pigmento incorporado ao composto antes da moldagem. A combinação dos três define o SKU (\textit{Stock Keeping Unit}, unidade de manutenção de estoque) sobre um produto de baixo valor unitário e alto volume, planejado e faturado em quilogramas de látex transformado.

A demanda por esses SKU é sazonal de calendário. O consumo se concentra em eventos datados, entre eles carnaval, festas juninas, formaturas, dia das crianças e festas de fim de ano, aos quais se somam campanhas de marca que exigem cores específicas em janelas curtas. A série de demanda é dominada por componente determinística de calendário, e não por um nível estacionário com ruído, o que desloca o esforço de previsão do patamar agregado para a datação e a composição de cada pico.

Esse padrão produz dois efeitos que se resolvem por decisões distintas. No agregado, a demanda alterna janelas que excedem a capacidade instalada e janelas que a subutilizam. No \textit{mix}, a composição do volume entre balões e entre cores muda de um evento para outro, de modo que um mesmo total pode corresponder a um conjunto inteiramente diferente de SKU. O primeiro efeito é de capacidade, o segundo é de alocação de máquina, e nenhum dos dois se resolve pela decisão que resolve o outro.

A capacidade não acompanha nenhum dos dois. O parque de máquinas de moldagem é fixo e o calendário de turnos é praticamente invariante ao longo do ano, porque dimensionar a fábrica para o pico deixaria o investimento ocioso na maior parte do horizonte. Atender o pico exige antecipar produção e carregar estoque sazonal, o que transfere a folga da capacidade para o estoque e o torna variável de decisão em vez de consequência do plano.

A tecnologia de processo limita o alcance dessa antecipação. A moldagem opera em fluxo contínuo a partir de um composto fluido mantido em tanques, e o ponto de diferenciação do produto está na pigmentação: até ali o composto é genérico, e depois dela está comprometido com uma cor. Trocar a cor em uma máquina exige limpeza do circuito e trocar o formato exige remontar o conjunto de formas, de modo que toda mudança de item consome tempo de \textit{setup}, horas de capacidade que não geram produto.

É o tempo de \textit{setup} que impede seguir a demanda SKU a SKU e impõe a produção em lotes, sob o compromisso corrente entre capacidade consumida em preparação e capital imobilizado em estoque. O que particulariza o caso dos balões é a estrutura desse \textit{setup}: a troca de cor é dependente de sequência e direcional, enquanto a troca de forma é uma intervenção mecânica de ordem de grandeza superior, decidida em cadência mais lenta. Com \textit{setup} dependente de sequência, dimensionar o lote deixa de ser separável de sequenciá-lo, e o problema é de dimensionamento e sequenciamento de lotes (\textit{lot sizing and scheduling}).

A demanda que alimenta essa decisão não chega por uma via única. Parte está confirmada em pedidos com quantidade, cor e prazo acordados, tratada aqui como carteira, e o restante é previsão. A operação é híbrida entre produção contra pedido (\textit{make-to-order}) e produção para estoque (\textit{make-to-stock}) sobre o mesmo parque de máquinas, com as duas correntes disputando a mesma capacidade. A distinção não é nominal: a falta contra carteira é atraso, transportado adiante e ainda devido, enquanto a falta contra previsão é venda perdida, que não se recupera.

Essa assimetria governa também o dimensionamento do estoque. O nível adequado por balão não é uma constante, e sim a soma de uma parcela de ciclo, imposta pelo tamanho de lote, de uma parcela sazonal, que antecipa o pico, e de uma parcela de cobertura, que absorve o erro de previsão no intervalo entre duas corridas do mesmo balão. Dimensioná-las é arbitrar entre o custo de carregamento, proporcional ao capital imobilizado, e o custo da falta, medido em margem perdida e em prazo descumprido.

As decisões resultantes não compartilham a mesma cadência. Quanto produzir de cada balão, em que máquina e em que semana é decisão revisada semanalmente, sobre um horizonte que precisa conter a subida sazonal; em que turno e em que ordem as cores entram em cada máquina é decisão da semana corrente, condicionada à carteira já existente. Separar essas duas cadências, e definir o que a primeira entrega à segunda, é o problema de planejamento formalizado na Seção~\ref{sec:modelo_semanal} e na Seção~\ref{sec:modelo_turno}.

\newpage

\section{Revisão de literatura}

\subsection{Planejamento hierárquico de produção}

O planejamento de produção organiza-se em níveis de decisão com horizontes e granularidades distintos: estratégico, tático e operacional \cite{anthony1965}. Decisões de nível superior fixam metas e restrições para o nível inferior; o nível inferior devolve ao superior o que de fato foi executado, fechando o ciclo.

\citeonline{haxmeal1975} formalizam esse esquema para produção multi-planta sob demanda sazonal: um plano agregado de estoque por família de produtos é obtido por programação linear e desagregado, em seguida, em programações detalhadas por item através de regras clássicas de controle de estoque. \citeonline{bitranhax1977} generalizam a proposta: o problema é particionado em subproblemas por nível, cada um resolvido por um procedimento próprio, com mecanismos explícitos de agregação (de baixo para cima) e desagregação (de cima para baixo) conectando os níveis.

Esse padrão é a base dos sistemas avançados de planejamento (\textit{Advanced Planning Systems}); \citeonline{stadtler2005} descreve a \textit{Supply Chain Planning Matrix}, que posiciona módulos de planejamento por horizonte (longo, médio e curto prazo) e por função (produção, distribuição e vendas), sendo a referência corrente para separar decisão tática de decisão operacional.

O mecanismo que liga os níveis é de duas vias. \citeonline{bitranhaashax1981,bitranhaashax1982} formalizam a desagregação em horizonte rolante: o plano agregado é resolvido para todo o horizonte, a primeira fatia é desagregada em famílias e itens sob restrições de consistência que obrigam a soma dos itens a reproduzir a meta agregada, e o realizado retorna ao nível superior na execução seguinte. \citeonline{schneeweiss2003} nomeia a contrapartida teórica desse arranjo: o nível superior só decide bem se carregar uma \textit{anticipation function}, um modelo aproximado da reação do nível inferior às instruções que recebe; sem ela, o plano agregado propaga inviabilidade para baixo em vez de restringi-la.

A correspondência família$\to$item de Hax e Meal mapeia diretamente o problema dos balões: formato e acabamento definem a família decidida no nível tático; a cor é o item desagregado no nível operacional. Essa fronteira não é uma escolha arbitrária de modelagem, e sim o reflexo da separação real entre os dois tipos de \textit{setup} do processo produtivo: a troca de forma, custosa e decidida com semanas de antecedência, e a troca de cor, uma limpeza decidida dia a dia.

\subsection{Dimensionamento de lotes capacitado}

O dimensionamento de lotes (\textit{lot sizing}) tem origem no modelo de \citeonline{wagnerwhitin1958}, que resolve por programação dinâmica o tamanho de lote ótimo de um único item sob demanda determinística variável no tempo, sem restrição de capacidade. A versão capacitada, chamada \textit{Capacitated Lot Sizing Problem} (CLSP), introduz um limite de capacidade por período e já é NP-difícil no caso de dois itens com \textit{setups} independentes \cite{bitranyanasse1982,chenthizy1990}.

O CLSP é dito de \textit{big-bucket}: cada período (tipicamente uma semana ou um mês) admite a produção de mais de um item, e a granularidade interna ao período não é modelada. \citeonline{karimi2003} revisam formulações e algoritmos exatos e heurísticos para essa classe. \citeonline{jansdegraeve2008} cobrem extensões que aproximam o modelo de aplicações industriais: máquinas paralelas, múltiplos níveis de estrutura de produto e \textit{setups} dependentes de sequência. \citeonline{buschkuhl2010} classificam quatro décadas de abordagens de solução. Para o caso de item único, \citeonline{brahimi2006,brahimi2017} cobrem formulações e métodos de solução em dois levantamentos que, juntos, resumem a literatura até 2016.

Aplicações industriais acrescentam ao CLSP canônico restrições que a formulação de livro-texto não tem. \citeonline{clark2011} mapeiam essas extensões e argumentam que os modelos acadêmicos padrão precisam ser reformulados caso a caso. Uma delas é o lote mínimo: \citeonline{andersoncheah1993} introduzem o tamanho mínimo de batelada junto com tempos de \textit{setup} num CLSP multi-item e mostram que a combinação exige decomposição para ser resolvida; \citeonline{okhrinrichter2011} caracterizam a versão em que o mínimo é imposto por quantidade mínima de pedido, sem custo fixo de \textit{setup}. Nos dois casos a forma padrão é $q_i\,y_{it} \le x_{it} \le M_{it}\,y_{it}$, com $y_{it}$ binária de \textit{setup}, e a qualidade da relaxação linear depende de $M_{it}$ ser o menor limitante válido disponível.

O nível semanal do problema dos balões é um CLSP em máquinas paralelas não relacionadas, com produtividade $p_{ij}$ distinta por par máquina-balão e compatibilidades próprias, ao qual se somam lote mínimo, custo de estoque explícito e duas camadas de demanda arbitradas por preço. É essa combinação que distingue o problema do CLSP canônico.

\subsection{Dimensionamento e sequenciamento integrados}

Quando o \textit{setup} depende da sequência de produção, como ocorre na troca de cor, o \textit{big-bucket} deixa de bastar: sem ordem definida dentro do período, não há como cobrar o \textit{setup} certo. A literatura resolve isso fundindo tamanho de lote e sequenciamento num único modelo de \textit{small-bucket}, com microperíodos dentro de cada período maior. O \textit{General Lotsizing and Scheduling Problem} (GLSP), de \citeonline{fleischmannmeyr1997}, é uma das formulações canônicas dessa família. O \textit{Capacitated Lot-sizing and Scheduling Problem with Sequence-Dependent setups} (CLSD), de \citeonline{haase1996}, é a outra. \citeonline{drexlkimms1997} organizam a taxonomia completa de modelos \textit{small-bucket}: o \textit{Discrete Lotsizing and Scheduling Problem} (DLSP), o \textit{Continuous Setup Lotsizing Problem} (CSLP), o \textit{Proportional Lotsizing and Scheduling Problem} (PLSP), além do GLSP e do CLSD já apresentados.

Um mecanismo central dessa literatura é o \textit{setup carryover}: a configuração da máquina persiste entre períodos consecutivos sem pagar novo \textit{setup}, desde que o mesmo item seja produzido em ambos. \citeonline{soxgao1999} introduzem o conceito para o CLSP; \citeonline{suriestadtler2003} formalizam a versão com lotes ligados (\textit{linked lot sizes}), permitindo que um lote atravesse a fronteira de dois períodos sem interrupção. \citeonline{copil2017} revisam e classificam toda essa família de modelos simultâneos de dimensionamento e sequenciamento; \citeonline{guimaraes2014} tratam especificamente da modelagem de \textit{setups} dependentes de sequência dentro desses modelos.

A escolha dentro dessa família não é indiferente. O DLSP, formalizado por \citeonline{fleischmann1990}, impõe a hipótese \textit{all-or-nothing}: em cada microperíodo ou se produz um único item ocupando toda a capacidade, ou não se produz nada. O PLSP de \citeonline{drexlhaase1995} relaxa essa hipótese admitindo no máximo dois itens por microperíodo. O GLSP não impõe nenhuma das duas: os microperíodos têm comprimento variável dentro de um macroperíodo de comprimento fixo, e o lote é contínuo. \citeonline{meyr2000,meyr2002} estendem o GLSP para \textit{setups} dependentes de sequência e para máquinas paralelas, que é a configuração da moldagem de balões.

O outro eixo de escolha é como cobrar dois \textit{setups} de escalas diferentes no mesmo modelo. \citeonline{gunther2014} formaliza o \textit{block planning} nesse papel: os produtos são agrupados em blocos, o \textit{setup} maior é pago apenas na fronteira do bloco e, dentro dele, a sequência percorre as variantes pagando somente o \textit{setup} menor. A economia de binárias vem de fixar a família dentro do bloco em vez de deixar a sequência livre sobre todas as combinações.

A troca de cor no problema dos balões é um \textit{setup} dependente de sequência e, além disso, direcional: a limpeza de claro para escuro é barata; de escuro para claro, cara. O \textit{carryover} de estado, isto é, a forma e a cor montadas na máquina, atravessa turnos e mesmo turnos ociosos, e é o mecanismo desta literatura reaproveitado no nível operacional do modelo proposto. \citeonline{menezes2011} tratam o caso vizinho em que a troca é longa demais para caber num único período, exigindo que o \textit{setup} atravesse a fronteira (\textit{crossover}); no problema dos balões a troca de forma cabe num turno, o que mantém o modelo no regime de \textit{carryover} puro.

\subsection{Demanda firme e demanda prevista}

A distinção entre carteira e previsão, estabelecida na Seção~\ref{sec:contextualizacao}, tem consequência direta sobre a formulação. O CLSP canônico trata a demanda de um item como corrente única, o que apaga a diferença entre uma falta de consequência contratual e uma falta que é venda não realizada, justamente onde essa diferença decide a alocação de capacidade escassa.

\citeonline{soman2004} organizam o problema sob o rótulo de produção híbrida \textit{make-to-order}/\textit{make-to-stock}: parte do mix é produzida contra pedido, parte contra previsão, e a decisão de quais itens vão para cada regime antecede o planejamento de capacidade. \citeonline{meyr2009} formaliza a consequência para o nível de serviço: quando a capacidade não cobre toda a demanda, a quantidade disponível é alocada entre classes de cliente com prioridades distintas, e a prioridade se expressa como diferença de penalidade, não como restrição rígida. \citeonline{slotnick2011} revisa a literatura de aceitação de pedidos, onde a mesma pergunta aparece invertida, a de qual pedido recusar ou atrasar quando não cabe tudo, e cataloga as formas de penalização usadas.

A mecânica de modelagem correspondente já existe na literatura de \textit{lot sizing}: \citeonline{belofilho2014} formulam o CLSP com \textit{backlogging} combinado a carryover e \textit{crossover} de \textit{setup}, mostrando como a falta transportada adiante convive com a preservação do estado da máquina entre períodos. O que não se encontra pronto é a combinação das duas naturezas de falta no mesmo modelo, a carteira transportada como atraso e a previsão perdida de vez, com penalidades assimétricas; essa é a forma adotada neste trabalho.

\subsection{Custo de manter estoque}
\label{sec:custo_estoque}

O custo de manter estoque não é um parâmetro de conveniência: é a contrapartida financeira do capital imobilizado e dos recursos consumidos para guardar o produto. \citeonline{lalondelambert1975} decompõem esse custo em quatro componentes, que são o custo de capital, o custo de serviço (seguro e impostos), o custo de espaço de armazenagem e o custo de risco (obsolescência, avaria, furto), e propõem a metodologia pela qual uma empresa calcula a própria taxa em vez de adotar um número de mercado.

A ordem de grandeza dessa taxa é estável na literatura. \citeonline{stocklambert2001} consolidam a regra prática de aproximadamente 25\,\% do valor do estoque ao ano. \citeonline{richardson1995} reporta uma faixa mais larga, de 25\,\% a 55\,\%, com a decomposição por faixa que explica a dispersão: o custo do dinheiro responde por 6\,\% a 12\,\% e a obsolescência por outros 6\,\% a 12\,\%, de modo que setores com giro lento ou produto perecível se deslocam para o topo da faixa. \citeonline{timme2003} argumentam que o componente de capital é sistematicamente subestimado e defendem o custo médio ponderado de capital (\textit{Weighted Average Cost of Capital}, WACC) como sua medida correta, com a partição típica de cerca de 15\,\% de capital e 10\,\% de custos não financeiros.

A forma operacional é a mesma em toda essa literatura: \citeonline{silver1998} definem o custo de manter uma unidade por unidade de tempo como $h = r\,v$, com $r$ a taxa de carregamento anual e $v$ o valor unitário do item, entendido como o custo variável acumulado até o estágio onde o estoque é mantido e não como o preço de venda, porque é esse o capital efetivamente parado. Adotar $r$ como percentual fixo do valor não é isento de crítica: \citeonline{berling2008} mostra, por custeio baseado em atividades, situações em que essa aproximação eleva o custo real em mais de 15\,\%, e \citeonline{azzi2014} documentam, em dez empresas, o quanto a taxa varia com o sistema de armazenagem empregado. \citeonline{castellanoglock2021} tratam do passo seguinte, que é o que interessa a um modelo de \textit{lot sizing}: como converter a taxa anual para o \textit{bucket} do modelo sem introduzir inconsistência na formulação de custo médio.

Neste trabalho a conversão é linear, $h_i = (\theta/52)\,v_i$, com $\theta$ a taxa anual e o estoque cobrado sobre a média entre o saldo de abertura e o de fechamento da semana. A taxa $\theta$ é parâmetro de entrada, declarado e sujeito a análise de sensibilidade, não constante embutida.

\subsection{Aplicações em indústria de processo}

Uma via alternativa para o dimensionamento e sequenciamento integrados evita inflar um único MILP (\textit{Mixed-Integer Linear Programming}, programação linear inteira mista): decompõe o problema em dois módulos, um de quantidade e outro de sequência, ou resolve a sequência por regra e não por otimização exata. Essa abordagem é comum em indústrias de processo com estrutura de dois estágios, onde um insumo intermediário é preparado numa etapa e diferenciado por sabor, cor ou aroma numa etapa seguinte.

A indústria de bebidas é a aplicação mais estudada e mais próxima estruturalmente do problema dos balões: xarope preparado num tanque, envasado em seguida em diferentes sabores e tamanhos. \citeonline{ferreira2009,ferreira2010} modelam esse problema de dois estágios, o preparo do xarope e o envase, com \textit{setups} dependentes de sequência e de tempo, comparando estratégias de \textit{relax-and-fix}; \citeonline{toledo2012} resolvem a mesma classe de problema com algoritmos meméticos. A analogia com os balões é direta: composto pigmentado é o xarope, moldagem é o envase.

A indústria têxtil de fiação oferece uma analogia ainda mais próxima, porque acrescenta uma restrição de sincronização que a de bebidas não tem. \citeonline{camargo2014} tratam um problema de dois estágios em que uma mistura de fibras (\textit{fiber blend}) alimenta múltiplas máquinas de fiação em paralelo; na formulação MSGLSP (\textit{Multi-Stage General Lot-sizing and Scheduling Problem}) que adotam, de \citeonline{camargo2012}, todas as máquinas ativas num mesmo microperíodo devem processar fios da mesma família de mistura, sincronizadas por uma variável binária de qualidade. Essa restrição, de haver só uma mistura disponível por vez e de todas as máquinas do segundo estágio precisarem respeitá-la, é exatamente o que emergiria se a pigmentação do composto de látex fosse modelada como etapa própria, com múltiplas moldagens em paralelo consumindo o mesmo composto pigmentado corrente. Duas diferenças, porém, mantêm o problema dos balões fora dessa formulação: os autores tratam a troca de mistura como \textit{setup} nulo (o oposto da troca de forma dos balões, que é o \textit{setup} caro) e adotam atraso (\textit{backorder}) com custo de estoque explícito, não venda perdida por custo de oportunidade.

Dois estudos de caso publicados em embalagens são os análogos mais próximos do ponto de vista do chão de fábrica. \citeonline{marinelli2007} tratam uma empresa de embalagens com máquinas paralelas e \textit{buffers} compartilhados, dimensionando lotes e sequenciando sob capacidade; \citeonline{almadalobo2007} formulam o caso da embalagem de vidro com \textit{setups} dependentes de máquina e de sequência, apresentando duas formulações com carryover de \textit{setup}, uma compacta e outra mais forte. Nenhum dos dois separa carteira de previsão nem cobra custo de estoque: a aproximação é de estrutura produtiva, não de estrutura de demanda.

\citeonline{gunther2006} descrevem o \textit{block planning}, aplicado originalmente à indústria de tintura de cabelo e depois generalizado para produção do tipo \textit{make-and-pack}: a sequência de produção segue uma ordem fixa e pré-definida de variantes, tipicamente por contaminação crescente, do tom mais claro para o mais escuro, evitando limpeza entre trocas. É a mesma lógica de sequenciamento por cor da moldagem de balões, com o mesmo argumento de fundo: claro antes de escuro minimiza limpeza.

\newpage

\section{Processo produtivo}

O processo produtivo (Figura \ref{fig:fluxo_produtivo}) inicia-se no recebimento de matérias-primas. Devido à sua perecibilidade, o látex é recebido diariamente e estocado em tanques, enquanto os demais insumos químicos pertencentes à produção são repostos conforme a demanda.

A transformação começa na pré-vulcanização, onde o látex é misturado a agentes químicos para viabilizar a futura elasticidade do produto. A seguir, o composto recebe pigmentos e aditivos de acabamento (perolado, brilhante, fosco), gerando diversas combinações demandadas pelo mercado.

Dimensionado pelo planejamento tático e sequenciado pelo planejamento operacional, o composto alimenta as máquinas de moldagem. Em um processo contínuo, os moldes imergem no fluido, recebem talco, são vulcanizados em estufas e desmoldados mecanicamente.

Por fim, batedores removem o excesso de talco dos balões semiacabados. Os produtos são então empacotados e transferidos para o estoque de produtos acabados, prontos para distribuição.

\begin{figure}[htbp]
    \centering
    \resizebox{\textwidth}{!}{% Redimensiona o diagrama para a largura do texto
        \begin{tikzpicture}[
            % Estilos do fluxograma
            process/.style={
                rectangle, 
                draw=black, 
                thick, 
                fill=white,
                text centered, 
                minimum height=1.8cm, 
                text width=3.2cm
            },
            inventory/.style={
                isosceles triangle,
                shape border rotate=180, % Inverte o triângulo
                draw=black,
                thick,
                fill=gray!25,
                text centered,
                minimum height=2.5cm,
                text width=2.5cm,
                inner ysep=0pt,
                text depth=1.5em % Ajusta o texto verticalmente no triângulo
            },
            critical/.style={
                process, % Herda o estilo de processo
                very thick,
                fill=gray!25,
                font=\bfseries
            },
            % Estilo das setas
            arrow/.style={
                draw=black, 
                -{>}, % Seta padrão do TikZ
                thick
            }
        ]
    
        % Nós do Processo
        % Linha Superior
        \node[process] (recebimento) {Recebimento de Matéria-Prima};
        \node[inventory, right=1.5cm of recebimento] (estoque_mp) {Estoque de Látex};
        \node[process, right=1.5cm of estoque_mp] (pre_vulcanizacao) {Pré-Vulcanização};
        \node[process, right=1.5cm of pre_vulcanizacao] (mistura_cor) {Mistura e Pigmentação};
    
        % Linha Inferior
        \node[critical, below=2cm of mistura_cor] (moldagem) {Moldagem, Vulcanização e Desmoldagem};
        \node[process, left=1.5cm of moldagem] (remocao_talco) {Remoção de Talco};
        \node[process, left=1.5cm of remocao_talco] (embalagem) {Embalagem};
        \node[inventory, left=1.5cm of embalagem] (estoque_pa) {Estoque de P.A.};
    
        % Conexões
        \path[arrow] (recebimento) -- (estoque_mp);
        \path[arrow] (estoque_mp) -- (pre_vulcanizacao);
        \path[arrow] (pre_vulcanizacao) -- (mistura_cor);
        \path[arrow] (mistura_cor) -- (moldagem);
        \path[arrow] (moldagem) -- (remocao_talco);
        \path[arrow] (remocao_talco) -- (embalagem);
        \path[arrow] (embalagem) -- (estoque_pa);
    
        \end{tikzpicture}%
    }
    \caption{Fluxograma do processo produtivo}
    \label{fig:fluxo_produtivo}
\end{figure}


\subsection{Planejamento tático}

O planejamento tático determina o que, onde, quando e quanto produzir, balanceando estoques e garantindo o atendimento da demanda ao longo do horizonte. A demanda e o estoque são agregados por formato e acabamento dos balões; a segregação por cor, granularidade inferior, não é considerada nesta etapa.

O nível de estoque alvo de cada período é definido em função da demanda futura: o tomador de decisão estabelece quantos períodos à frente o estoque deve cobrir, com o alvo (de cada combinação de formato e acabamento de balão) assumindo valores que podem variar dinamicamente conforme a demanda (prevista ou passada). Essa cobertura orienta se a produção do período deve repor ou consumir o estoque disponível.

O plano de produção é elaborado em periodicidade semanal, com revisões eventuais conforme a necessidade operacional. As trocas de formas entre balões de diferentes formatos configuram longos tempos de \textit{setup}: no limite mínimo, ocorrem entre turnos de trabalho, mas na prática se dão com intervalos de dias, dado o alto custo de tempo envolvido. A heterogeneidade entre equipamentos é o fator que viabiliza a operação simultânea de diferentes balões: cada máquina pode ser alocada a um formato distinto, permitindo atender o \textit{mix} sem que todas as máquinas passem por \textit{setup} ao mesmo tempo.

Cada máquina possui taxas de produção e compatibilidades próprias para cada formato e acabamento de balão. A capacidade produtiva pode ser reduzida por ociosidades programadas, decorrentes de manutenções preventivas, reformas ou paradas não planejadas.

A proporção entre carteira e previsão varia ao longo do horizonte: nas semanas próximas quase toda a demanda já está vendida, e à medida que o horizonte se afasta a carteira rareia e a previsão predomina. O planejamento tático serve as duas correntes numa mesma decisão de capacidade, sabendo que falhar com um pedido firme e falhar com uma estimativa não têm a mesma consequência.

O objetivo do modelo é minimizar o custo total da decisão, em reais: a margem perdida nas faltas, ponderada pela natureza de cada uma, o custo direto das horas de \textit{setup} e o custo de carregar estoque. Essa estrutura é formalizada na função objetivo do modelo semanal, apresentado na Seção~\ref{sec:modelo_semanal}.

\subsection{Previsão de demanda}

A previsão da demanda é um elemento chave do planejamento, pois é a partir dela que se definem a alocação das máquinas, e por consequência as decisões de estocagem e, possivelmente, priorização de perdas por não atendimento de demanda.

No contexto do planejamento tático no problema de produção de balões de látex, a acurácia da previsão é fundamental e depende fortemente de conseguir capturar períodos de sazonalidade ao longo do ano, bem como tendências da empresa e do mercado.

Nessa etapa da modelagem do problema, o objetivo do estudo é explorar métodos e arquiteturas diferentes de previsão de demanda a fim de observar o impacto da acuracidade no planejamento produtivo. O parâmetro $\alpha$ no modelo de otimização, por sua vez, pode ser usado para controlar o nível de estoque e reduzir o impacto do erro dos modelos: melhores acuracidades permitem maior giro de estoque ao longo do ano.


\subsection{Planejamento operacional}

O planejamento operacional detalha a primeira semana do plano tático. Enquanto o nível tático decide quanto produzir de cada balão, o operacional decide em que turno, em que máquina e em que ordem; é nele que a cor aparece, porque é a cor que determina a limpeza entre uma corrida e a seguinte.

A unidade de decisão é o turno. A fábrica opera em três turnos diários, e a troca de forma consome parcela relevante de um turno de oito horas: na prática ela é feita na virada, quando a máquina já está parando. A troca de cor é mais curta e ocorre dentro do turno, entre corridas do mesmo balão. Essa diferença de escala é o que separa as duas decisões: a forma montada numa máquina vale para o turno inteiro, e o que se sequencia dentro dele são as cores.

A troca de cor é direcional e nem sempre possível. Passar de um tom claro para um escuro exige pouca limpeza; o caminho inverso exige muito mais, e alguns pares simplesmente não são executados na sequência imediata, por risco de contaminação do composto. A fábrica mantém essa informação como uma matriz de tempos entre pares de cor, com células vazias marcando as transições proibidas.

Uma ordem de produção não atravessa a virada de turno: o que é iniciado num turno termina nele. Some-se a isso o lote mínimo, abaixo do qual a corrida não compensa o tempo de preparação, e o conjunto de decisões do nível operacional fica definido: quais operações cada máquina executa em cada turno, em que sequência, e em que quantidade.

\subsection{Granularidade das decisões}

As decisões se distribuem por três escalas de tempo encaixadas, e a escolha de cada uma responde a um fenômeno físico do processo, não a conveniência de modelagem.

A \textbf{semana} é a escala do plano tático. É o intervalo em que a fábrica revisa o plano e em que a demanda é comercialmente estável o bastante para ser tratada como dado. O horizonte cobre um número de semanas suficiente para conter a subida sazonal, porque antecipar produção contra o pico é uma decisão que só existe se o modelo enxerga o pico.

O \textbf{turno} é a escala do plano operacional. É a menor unidade em que a fábrica organiza mão de obra e a fronteira natural das trocas de forma. Só a primeira semana do plano tático é detalhada em turnos; as demais permanecem agregadas até a revisão seguinte, num esquema de horizonte rolante.

A \textbf{posição dentro do turno} é a escala da sequência. Cada turno comporta um número pequeno de corridas de cor, porque o lote mínimo e o tempo de limpeza limitam quantas cabem, e a ordem entre elas é o que determina qual \textit{setup} de cor é pago. Não há escala abaixo dessa: o modelo não decide instantes, decide a ordem.

O que desce da semana para o turno é uma única grandeza: o estoque que cada balão deve ter ao fim da semana. O que sobe de volta, na revisão seguinte, é o estado real: estoque, carteira consumida e a configuração em que cada máquina ficou.

\newpage

\section{Modelo de demanda}
\label{sec:modelo_demanda}

A modelagem da demanda proposta neste trabalho baseia-se em uma abordagem econométrica hierárquica \textit{top-down}. O objetivo é isolar os efeitos exógenos e sazonais em níveis agregados de histórico de vendas, onde o sinal estatístico é mais forte e o ruído é menor, e repassar esse aprendizado para os níveis mais granulares na forma de coeficientes \textit{a priori}. 

O modelo extrai características (\textit{features}) exógenas e temporais através de múltiplos níveis de agregação e, em seguida, normaliza as vendas para a obtenção de uma demanda base por produto. O \textit{forecast} final é a recombinação dessa base com os efeitos sazonais projetados.

\subsection*{Conjuntos}
\begin{tabularx}{\linewidth}{@{}L{2.8cm}R@{}}
$T$ & Conjunto de períodos históricos.\\
$I$ & Conjunto de balões (combinação de formato e acabamento).\\
$F$ & Conjunto de formatos de balões.\\
$I_f \subseteq I$ & Subconjunto de balões pertencentes ao formato $f \in F$.\\
$K$ & Conjunto de \textit{features} exógenas e sazonais.\\
\end{tabularx}

\subsection*{Parâmetros}
\begin{tabularx}{\linewidth}{@{}L{2.8cm}R@{}}
$y_{it}$ & Demanda real observada do balão $i$ no período $t$.\\
$Y_{ft}$ & Demanda real agregada do formato $f$ no período $t$.\\
$Y_{t}$ & Demanda real global agregada no período $t$.\\
$X_{kt}$ & Valor da \textit{feature} $k$ no período $t$.\\
$\lambda$ & Fator de penalização na transferência de coeficientes.\\
$h$ & Horizonte futuro de previsão.\\
\end{tabularx}

\subsection*{Variáveis}
\begin{tabularx}{\linewidth}{@{}L{2.8cm}R@{}}
$\beta^{(0)}_{k}$ & Coeficiente da \textit{feature} $k$ no nível global a ser estimado.\\
$\beta^{(1)}_{fk}$ & Coeficiente da \textit{feature} $k$ para o formato $f$ a ser estimado.\\
$\beta^{(2)}_{ik}$ & Coeficiente da \textit{feature} $k$ para o balão $i$ a ser estimado.\\
$\hat{\beta}$ & Parâmetros já estimados e fixados na etapa anterior de agregação.\\
$S_{it}$ & Efeito sazonal/exógeno consolidado para o balão $i$ no período $t$.\\
$\tilde{y}_{it}$ & Demanda dessazonalizada do balão $i$ no período $t$.\\
$B_{i}$ & Demanda base média do balão $i$.\\
$\hat{d}_{i,t+h}$ & Previsão final de demanda do balão $i$ para o período futuro $t+h$.\\
\end{tabularx}

\vspace{1cm}

A extração dos efeitos sazonais e exógenos é feita em cascata. O modelo utiliza uma formulação log-linear (multiplicativa) em detrimento de uma abordagem aditiva devido à natureza da propagação da demanda. Em eventos de forte sazonalidade (final de ano, por exemplo), a demanda agregada pode triplicar. A transformação logarítmica permite que esse comportamento atue como um fator multiplicativo percentual, sendo repassado de forma coerente e proporcional para os níveis inferiores. Caso um modelo aditivo fosse empregado, um incremento absoluto capturado no nível global (um aumento sazonal de 10.000 kg, por exemplo) perderia o sentido ao ser cascateado para uma combinação específica de formato e acabamento, cuja escala natural de vendas é substancialmente menor. Ademais, a modelagem multiplicativa previne intrinsecamente a projeção de demandas negativas.

Para a estimação destes parâmetros, utiliza-se o algoritmo \textit{Recursive Least Squares} (RLS), que atualiza os coeficientes incrementalmente a cada nova observação minimizando o erro quadrático acumulado.

A estimação inicial, no nível global, ocorre sobre a demanda total consolidada:
\begin{ceqn}
\begin{equation}
\hat{\beta}^{(0)} = \arg\min_{\beta} \sum_{t \in T} \left( \ln(Y_t) - \sum_{k \in K} \beta_{k} X_{kt} \right)^2
\label{eq:demanda_nivel0}
\end{equation}
\end{ceqn}

No nível de formato, os coeficientes globais $\hat{\beta}^{(0)}_k$ atuam como \textit{priors} para a estimação, e o termo de regularização $\lambda$ penaliza desvios do comportamento global:
\begin{ceqn}
\begin{equation}
\hat{\beta}^{(1)}_{f} = \arg\min_{\beta} \left( \sum_{t \in T} \left( \ln(Y_{ft}) - \sum_{k \in K} \beta_{k} X_{kt} \right)^2 + \lambda \sum_{k \in K} \left( \beta_{k} - \hat{\beta}^{(0)}_k \right)^2 \right) \quad \forall f \in F
\label{eq:demanda_nivel1}
\end{equation}
\end{ceqn}

No nível de formato-acabamento, por fim, os coeficientes do formato respectivo atuam como \textit{priors} para o balão específico:
\begin{ceqn}
\begin{equation}
\hat{\beta}^{(2)}_{i} = \arg\min_{\beta} \left( \sum_{t \in T} \left( \ln(y_{it}) - \sum_{k \in K} \beta_{k} X_{kt} \right)^2 + \lambda \sum_{k \in K} \left( \beta_{k} - \hat{\beta}^{(1)}_{f,k} \right)^2 \right) \quad \forall f \in F, i \in I_f
\label{eq:demanda_nivel2}
\end{equation}
\end{ceqn}

Com os coeficientes granulares $\hat{\beta}^{(2)}_{ik}$ obtidos, o efeito consolidado para o balão $i$ no período $t$ retorna à escala original:
\begin{ceqn}
\begin{equation}
S_{it} = \exp \left( \sum_{k \in K} \hat{\beta}^{(2)}_{ik} X_{kt} \right) \quad \forall i \in I, t \in T
\label{eq:efeito_consolidado}
\end{equation}
\end{ceqn}

A demanda observada é então normalizada pelo efeito consolidado, isolando a demanda base do produto naquele período:
\begin{ceqn}
\begin{equation}
\tilde{y}_{it} = \frac{y_{it}}{S_{it}} \quad \forall i \in I, t \in T
\label{eq:demanda_normalizada}
\end{equation}
\end{ceqn}

A demanda base é definida como a média das vendas normalizadas ao longo de todo o horizonte histórico observado $|T|$:
\begin{ceqn}
\begin{equation}
B_{i} = \frac{1}{|T|} \sum_{t \in T} \tilde{y}_{it} \quad \forall i \in I
\label{eq:demanda_base}
\end{equation}
\end{ceqn}

O \textit{forecast} para um período futuro $t+h$ projeta a demanda base multiplicando-a pelos efeitos esperados:
\begin{ceqn}
\begin{equation}
\hat{d}_{i,t+h} = B_{i} \cdot \exp \left( \sum_{k \in K} \hat{\beta}^{(2)}_{ik} X_{k,t+h} \right) \quad \forall i \in I
\label{eq:forecast_final}
\end{equation}
\end{ceqn}

\newpage

\section{Modelo de otimização}
\label{sec:modelo_otimizacao}

O problema é resolvido em dois níveis encadeados. O nível semanal decide, para todo o horizonte de planejamento, quanto produzir de cada balão em cada máquina e em cada semana, arbitrando entre carteira e previsão. O nível operacional recebe desse plano uma única informação, a meta de estoque ao fim da primeira semana, e decide, para essa semana congelada, o que cada máquina produz em cada turno e em que ordem. As duas formulações compartilham a mesma unidade monetária e os mesmos parâmetros de custo; diferem no \textit{bucket} de tempo e no nível de agregação do produto.

\subsection{Notação comum}

Os índices são fixos nos dois modelos: $i$ para balão (formato e acabamento, sem cor), $j$ para máquina, $t$ para semana, $s$ para turno, $n$ para posição dentro do turno e $k$ para SKU-cor, isto é, a combinação de balão e cor. O SKU é sempre tratado junto com o balão a que pertence, por meio do subconjunto $K_i$ definido adiante, de modo que os parâmetros indexados por balão permanecem escritos em $i$.

\begin{tabularx}{\linewidth}{@{}L{2.8cm}R@{}}
$v_{i}$ & Valor unitário do balão $i$: custo variável de produção acumulado (R\$/kg).\\
$m_{i}$ & Margem de contribuição unitária do balão $i$: preço de venda menos $v_i$ (R\$/kg).\\
$c_{j}$ & Custo horário de \textit{setup} da máquina $j$: mão de obra, energia e composto descartado (R\$/h).\\
$p_{ij}$ & Produtividade do balão $i$ na máquina $j$ (kg/hora).\\
$q_{i}$ & Lote mínimo do balão $i$ (kg).\\
$\theta$ & Taxa anual de carregamento de estoque (Seção~\ref{sec:custo_estoque}).\\
\end{tabularx}

Os três parâmetros de custo são distintos por construção. A venda que deixa de acontecer custa a margem $m_i$, não o custo $v_i$; o estoque parado custa sobre $v_i$, que é o capital imobilizado; e o \textit{setup} custa as horas de máquina e equipe que consome, medidas por $c_j$. O custo de oportunidade das horas de \textit{setup} não entra na função objetivo: essas horas já reduzem a capacidade disponível na restrição correspondente, e a produção que deixa de ocorrer reaparece como falta penalizada; cobrá-la também no objetivo contaria a mesma perda duas vezes.

\subsection{Modelo semanal}
\label{sec:modelo_semanal}

O nível semanal é um problema de dimensionamento de lotes capacitado em máquinas paralelas não relacionadas. O índice de dia não existe neste nível: a ordem de produção dentro da semana é decisão do nível operacional, e o \textit{setup} é cobrado por ativação de par balão-máquina na semana, independente de sequência.

\subsubsection*{Conjuntos}
\begin{tabularx}{\linewidth}{@{}L{2.8cm}R@{}}
$T$ & Conjunto de semanas do horizonte de planejamento.\\
$I$ & Conjunto de balões.\\
$J$ & Conjunto de máquinas ativas.\\
$J_{i} \subseteq J$ & Máquinas compatíveis com o balão $i$.\\
\end{tabularx}

\subsubsection*{Parâmetros}
\begin{tabularx}{\linewidth}{@{}L{2.8cm}R@{}}
$o_{it}$ & Carteira firme do balão $i$ com vencimento na semana $t$ (kg).\\
$\hat{d}_{it}$ & Previsão líquida do balão $i$ na semana $t$: previsão total menos $o_{it}$, limitada inferiormente a zero (kg).\\
$\kappa_{jt}$ & Horas produtivas disponíveis da máquina $j$ na semana $t$.\\
$\tau_{i}$ & Tempo médio de \textit{setup} para o balão $i$ (horas).\\
$I^{0}_{i}$ & Estoque do balão $i$ no início do horizonte (kg).\\
\end{tabularx}

\subsubsection*{Variáveis}
\begin{tabularx}{\linewidth}{@{}L{2.8cm}R@{}}
$x_{ijt} \ge 0$ & Quantidade do balão $i$ produzida na máquina $j$ na semana $t$ (kg).\\
$y_{ijt} \in \{0,1\}$ & 1 se a máquina $j$ é preparada para o balão $i$ na semana $t$.\\
$I_{it} \ge 0$ & Estoque do balão $i$ ao fim da semana $t$ (kg).\\
$a_{it} \ge 0$ & Carteira do balão $i$ ainda não entregue ao fim da semana $t$ (kg).\\
$l_{it} \ge 0$ & Previsão do balão $i$ perdida na semana $t$ (kg).\\
\end{tabularx}

\subsection*{\textbf{Minimizar}}
\begin{ceqn}
\begin{equation}
\begin{aligned}
Z^{S} = {} & \sum_{i \in I}\sum_{t \in T} m_{i} \bigl( a_{it} + l_{it} \bigr)
    + \sum_{i \in I}\sum_{j \in J_i}\sum_{t \in T} c_{j}\, \tau_{i}\, y_{ijt} \\
    & + \sum_{i \in I}\sum_{t \in T} \frac{\theta}{52}\, v_{i}\, \frac{I_{i,t-1} + I_{it}}{2}
\end{aligned}
\label{eq:obj_semanal}
\end{equation}
\end{ceqn}

\noindent\textbf{Sujeito a}
\begin{align}
I_{i,t-1} + \sum_{j \in J_i} x_{ijt}
    &= I_{it} + \bigl(o_{it} + a_{i,t-1} - a_{it}\bigr) + \bigl(\hat{d}_{it} - l_{it}\bigr)
    && \forall\, i \in I,\; t \in T \label{eq:balanco_semanal} \\
a_{it} &\le o_{it} + a_{i,t-1}
    && \forall\, i \in I,\; t \in T \label{eq:atraso_limite} \\
l_{it} &\le \hat{d}_{it}
    && \forall\, i \in I,\; t \in T \label{eq:perda_limite} \\
\sum_{i \in I} \left( \frac{x_{ijt}}{p_{ij}} + \tau_{i}\, y_{ijt} \right)
    &\le \kappa_{jt}
    && \forall\, j \in J,\; t \in T \label{eq:capacidade_semanal} \\
q_{i}\, y_{ijt} \le x_{ijt} &\le M_{ijt}\, y_{ijt}
    && \forall\, i \in I,\; j \in J_i,\; t \in T \label{eq:lote_minimo}
\end{align}

A função objetivo \eqref{eq:obj_semanal} soma quatro parcelas, todas em reais. As duas primeiras são as faltas, ambas cobradas à margem $m_i$: o atraso da carteira e a previsão perdida. A terceira cobra as horas de \textit{setup} ao custo horário da máquina, e a quarta o carregamento de estoque sobre o saldo médio da semana, com a taxa anual $\theta$ rateada pelas 52 semanas do ano.

Cobrar as duas faltas ao mesmo preço unitário não as torna equivalentes. O atraso $a_{it}$ entra no somatório \emph{em toda semana em que persiste}, e a carteira atrasada permanece no balanço \eqref{eq:balanco_semanal}: ela ainda terá de ser produzida. A previsão perdida $l_{it}$ é cobrada uma única vez e sai do balanço, poupando a produção correspondente. Atrasar um quilograma por uma semana custa, portanto, $m_i$ mais o custo de produzi-lo adiante, ao passo que perdê-lo custa $m_i$ e dispensa a produção. A carteira é estritamente mais cara que a previsão por construção da formulação, sem que seja necessário um multiplicador de calibração. Nenhuma das duas faltas é proibida: a hierarquia entre elas decorre da estrutura do balanço, o que mantém o modelo viável sob qualquer capacidade e faz a prioridade emergir da otimização em vez de ser imposta por restrição ou por um peso arbitrado fora do modelo.

O balanço \eqref{eq:balanco_semanal} separa as duas naturezas de demanda no mesmo estoque. O termo $o_{it} + a_{i,t-1} - a_{it}$ é a carteira efetivamente entregue na semana, ou seja, o que venceu mais o que vinha atrasado menos o que segue atrasado, de modo que a falta de pedido firme não desaparece: é transportada adiante e volta a ser cobrada. O termo $\hat{d}_{it} - l_{it}$ é a previsão atendida, e o que não se atende em $l_{it}$ está perdido. A restrição \eqref{eq:atraso_limite} impede que o atraso exceda a carteira acumulada, o que tornaria a entrega negativa; \eqref{eq:perda_limite} faz o mesmo do lado da previsão.

Não há terceira camada de demanda. Uma meta de cobertura semana a semana, que exigisse estoque não inferior à demanda das semanas seguintes, seria redundante com o próprio horizonte: o modelo já enxerga essa demanda dentro de $T$ e antecipa produção por conta própria quando a capacidade aperta, comportamento característico do CLSP. A única função que tal meta cumpriria é conter o efeito de borda, em que o estoque é drenado a zero na última semana por não haver demanda além de $T$, e esse é um artefato de horizonte finito, tratado pelo horizonte rolante descrito nas premissas, em que apenas a primeira semana é congelada e o plano é refeito a cada execução. Acrescentar uma meta $\sigma_{it}$ e um peso $\rho$ para esse fim introduziria dois parâmetros sem contrapartida observável na operação, cuja calibração recairia sobre quem opera o modelo.

A capacidade \eqref{eq:capacidade_semanal} converte quantidade em horas pela produtividade e acrescenta as horas de \textit{setup} de cada par ativado, ambas competindo pelas $\kappa_{jt}$ horas disponíveis. O lote mínimo \eqref{eq:lote_minimo} liga quantidade e ativação nos dois sentidos: produzir exige preparar a máquina, e preparar exige produzir ao menos $q_i$. O limitante $M_{ijt}$ é tomado como o menor valor válido entre a capacidade da máquina na semana, $p_{ij}\,\kappa_{jt}$, e a demanda remanescente do balão no horizonte; um $M$ frouxo degradaria a relaxação linear e, com ela, o desempenho do \textit{branch-and-bound} \cite{jansdegraeve2008}.

\subsection{Modelo operacional por turno}
\label{sec:modelo_turno}

O nível operacional detalha a primeira semana do plano com \textit{bucket} de turno. A unidade de decisão é o par máquina-turno: nele uma única forma fica montada, e o turno é dividido em $\bar{n}$ posições ordenadas, que são os microperíodos em que a sequência de cores se realiza. É a estrutura de macroperíodo e microperíodos do GLSP \cite{fleischmannmeyr1997}, com o macroperíodo fixado no turno e a forma fixada no macroperíodo, o que corresponde à adaptação do \textit{block planning} de \citeonline{gunther2014} a este processo: o bloco é o turno, o \textit{setup} maior é pago na fronteira e a sequência interna percorre apenas as cores.

A notação acompanha a do nível semanal sempre que o papel do símbolo é o mesmo. As quantidades produzidas continuam em $x$, as binárias de alocação em $y$, o estoque em $I$, o atraso da carteira em $a$, a previsão perdida em $l$, a violação de meta em $g$ e as horas disponíveis em $\kappa$. O que é próprio deste nível são os dois tipos de \textit{setup}, e eles se distinguem por um sobrescrito: $\mathrm{F}$ para forma e $\mathrm{C}$ para cor. Assim, $w^{\mathrm{F}}_{ijs}$ e $w^{\mathrm{C}}_{kjs}$ são os dois estados de configuração da máquina, $z^{\mathrm{F}}$ e $z^{\mathrm{C}}$ são as trocas correspondentes, e $\tau^{\mathrm{F}}$ e $\tau^{\mathrm{C}}$ são seus tempos, na mesma letra $\tau$ que o nível semanal usa para o tempo médio de \textit{setup}.

O produto aparece em dois níveis de agregação. O balão $i$ continua sendo a unidade da forma, do lote mínimo e da produtividade; o SKU $k$ acrescenta a cor e é a unidade da demanda, do estoque e da sequência. O conjunto $K_i$ reúne os SKU do balão $i$ e particiona $K$, de modo que um somatório duplo sobre $i \in I$ e $k \in K_i$ percorre todos os SKU e mantém disponível o índice $i$ de que os parâmetros por balão dependem.

\subsubsection*{Conjuntos}
\begin{tabularx}{\linewidth}{@{}L{2.8cm}R@{}}
$S$ & Turnos da semana congelada, ordenados; $\bar{s}$ é o último.\\
$N = \{1,\dots,\bar{n}\}$ & Posições dentro de um turno, ordenadas.\\
$K$ & Conjunto de SKU, isto é, pares de balão e cor.\\
$K_{i} \subseteq K$ & SKU do balão $i$; os conjuntos $K_i$ particionam $K$.\\
$A \subseteq K \times K$ & Pares de SKU cuja transição imediata é permitida.\\
\end{tabularx}

\subsubsection*{Parâmetros}
\begin{tabularx}{\linewidth}{@{}L{2.8cm}R@{}}
$\kappa_{js}$ & Horas produtivas disponíveis da máquina $j$ no turno $s$.\\
$\tau^{\mathrm{F}}_{i'i}$ & Tempo de \textit{setup} para trocar a forma do balão $i'$ para a do balão $i$ (horas).\\
$\tau^{\mathrm{C}}_{kk'}$ & Tempo de \textit{setup} para trocar a cor de $k$ para $k'$, definido nos pares de $A$ (horas).\\
$o_{ks}$ & Carteira firme do SKU $k$ com vencimento no turno $s$ (kg).\\
$\hat{d}_{ks}$ & Previsão líquida do SKU $k$ atribuída ao turno $s$ (kg).\\
$\bar{I}_{i}$ & Estoque do balão $i$ ao fim da semana escolhido pelo plano semanal (kg).\\
$I^{0}_{k}$ & Estoque do SKU $k$ no início da semana (kg).\\
\end{tabularx}

\subsubsection*{Variáveis}
\begin{tabularx}{\linewidth}{@{}L{2.8cm}R@{}}
$x_{kjsn} \ge 0$ & Quantidade do SKU $k$ produzida na posição $n$ do turno $s$ na máquina $j$ (kg).\\
$y_{kjsn} \in \{0,1\}$ & 1 se o SKU $k$ ocupa essa posição.\\
$w^{\mathrm{F}}_{ijs} \in [0,1]$ & 1 se a máquina $j$ está com a forma do balão $i$ montada no turno $s$.\\
$w^{\mathrm{C}}_{kjs} \in [0,1]$ & 1 se a máquina $j$ termina o turno $s$ com a cor do SKU $k$ montada.\\
$u_{js} \in [0,1]$ & 1 se a máquina $j$ produz no turno $s$.\\
$z^{\mathrm{F}}_{i'ijs} \in [0,1]$ & 1 se a máquina $j$ troca a forma de $i'$ para $i$ na abertura do turno $s$.\\
$z^{\mathrm{C}}_{kk'jsn} \in [0,1]$ & 1 se a cor passa de $k$ para $k'$ ao entrar na posição $n$.\\
$I_{ks}, a_{ks}, l_{ks} \ge 0$ & Estoque, carteira em atraso e previsão perdida do SKU $k$ ao fim do turno $s$ (kg).\\
$e^{o}_{ks}, e^{d}_{ks} \ge 0$ & Quantidade do SKU $k$ entregue no turno $s$ contra carteira e contra previsão (kg).\\
$g_{i} \ge 0$ & Violação da meta de estoque do balão $i$ ao fim da semana (kg).\\
\end{tabularx}

Somente $y_{kjsn}$ é declarada binária. As demais variáveis de configuração são contínuas em $[0,1]$ e assumem valor inteiro na solução ótima por consequência das restrições: $w^{\mathrm{F}}$ e $u$ ficam presos ao valor de $y$ pelas equações \eqref{eq:posicao_ocupada} e \eqref{eq:forma_da_posicao}, enquanto $w^{\mathrm{C}}$, $z^{\mathrm{F}}$ e $z^{\mathrm{C}}$ são minimizados contra limitantes inferiores do tipo $\text{AND}$ linearizado. Declarar apenas $y$ como binária reduz substancialmente o número de variáveis inteiras sem relaxar o modelo.

\subsection*{\textbf{Minimizar}}
\begin{ceqn}
\begin{equation}
\begin{aligned}
Z^{O} = {} & \sum_{i \in I}\sum_{k \in K_i}\sum_{s \in S} m_{i} \bigl( a_{ks} + l_{ks} \bigr)
    + \sum_{i \in I} m_{i}\, g_{i} \\
    & + \sum_{j \in J}\sum_{s \in S} c_{j} \left(
        \sum_{i' \in I}\sum_{i \in I} \tau^{\mathrm{F}}_{i'i}\, z^{\mathrm{F}}_{i'ijs}
        + \sum_{(k,k') \in A}\sum_{n \in N} \tau^{\mathrm{C}}_{kk'}\, z^{\mathrm{C}}_{kk'jsn} \right) \\
    & + \sum_{i \in I}\sum_{k \in K_i}\sum_{s \in S} \frac{\theta}{52\,|S|}\, v_{i}\, \frac{I_{k,s-1} + I_{ks}}{2}
\end{aligned}
\label{eq:obj_turno}
\end{equation}
\end{ceqn}

\noindent\textbf{Sujeito a}
\begin{align}
\sum_{i \in I} w^{\mathrm{F}}_{ijs} &\le 1
    && \forall\, j \in J,\; s \in S \label{eq:forma_unica} \\
u_{js} &\le \sum_{i \in I} w^{\mathrm{F}}_{ijs}
    && \forall\, j \in J,\; s \in S \label{eq:maquina_ativa} \\
\sum_{k \in K} y_{kjsn} &= u_{js}
    && \forall\, j \in J,\; s \in S,\; n \in N \label{eq:posicao_ocupada} \\
y_{kjsn} &\le w^{\mathrm{F}}_{ijs}
    && \forall\, i \in I,\; k \in K_i,\; j \in J,\; s \in S,\; n \in N \label{eq:forma_da_posicao} \\
z^{\mathrm{F}}_{i'ijs} &\ge w^{\mathrm{F}}_{i'j,s-1} + w^{\mathrm{F}}_{ijs} - 1
    && \forall\, i' \neq i,\; j \in J,\; s \in S \label{eq:troca_forma} \\
z^{\mathrm{C}}_{kk'jsn} &\ge y_{kjs,n-1} + y_{k'jsn} - 1
    && \forall\, (k,k') \in A,\; j \in J,\; s \in S,\; n \in N \label{eq:troca_cor} \\
y_{kjs,n-1} + y_{k'jsn} &\le 1
    && \forall\, (k,k') \notin A,\; j \in J,\; s \in S,\; n \in N \label{eq:transicao_proibida} \\
\sum_{k \in K_i} w^{\mathrm{C}}_{kjs} &= w^{\mathrm{F}}_{ijs}
    && \forall\, i \in I,\; j \in J,\; s \in S \label{eq:estado_cor} \\
w^{\mathrm{C}}_{kjs} &\ge y_{kjs\bar{n}}
    && \forall\, k \in K,\; j \in J,\; s \in S \label{eq:estado_ultima} \\
\bigl| w^{\mathrm{F}}_{ijs} - w^{\mathrm{F}}_{ij,s-1} \bigr| &\le u_{js}
    && \forall\, i \in I,\; j \in J,\; s \in S \label{eq:carryover_forma} \\
\bigl| w^{\mathrm{C}}_{kjs} - w^{\mathrm{C}}_{kj,s-1} \bigr| &\le u_{js}
    && \forall\, k \in K,\; j \in J,\; s \in S \label{eq:carryover_cor}
\end{align}
\begin{ceqn}
\begin{equation}
\begin{aligned}
\sum_{i \in I}\sum_{k \in K_i}\sum_{n \in N} \frac{x_{kjsn}}{p_{ij}}
    + \sum_{i' \in I}\sum_{i \in I} \tau^{\mathrm{F}}_{i'i}\, z^{\mathrm{F}}_{i'ijs}
    + \sum_{(k,k') \in A}\sum_{n \in N} \tau^{\mathrm{C}}_{kk'}\, z^{\mathrm{C}}_{kk'jsn}
    \;\le\; \kappa_{js}
    \qquad \forall\, j \in J,\; s \in S
\end{aligned}
\label{eq:capacidade_turno}
\end{equation}
\end{ceqn}

\begin{ceqn}
\begin{equation}
\begin{aligned}
q_{i} \bigl( y_{kjsn} - y_{kjs,n-1} \bigr) \;\le\; x_{kjsn} \;\le\; p_{ij}\, \kappa_{js}\, y_{kjsn}
    \qquad \forall\, i \in I,\; k \in K_i,\; j \in J,\; s \in S,\; n \in N
\end{aligned}
\label{eq:lote_minimo_turno}
\end{equation}
\end{ceqn}

\begin{align}
I_{ks} &= I_{k,s-1} + \sum_{j \in J}\sum_{n \in N} x_{kjsn} - e^{o}_{ks} - e^{d}_{ks}
    && \forall\, k \in K,\; s \in S \label{eq:balanco_turno} \\
a_{ks} &= a_{k,s-1} + o_{ks} - e^{o}_{ks}
    && \forall\, k \in K,\; s \in S \label{eq:atraso_turno} \\
e^{d}_{ks} + l_{ks} &= \hat{d}_{ks}
    && \forall\, k \in K,\; s \in S \label{eq:previsao_turno} \\
\sum_{k \in K_i} I_{k,\bar{s}} + g_{i} &\ge \bar{I}_{i}
    && \forall\, i \in I \label{eq:meta_final}
\end{align}

A função objetivo \eqref{eq:obj_turno} tem as mesmas parcelas de \eqref{eq:obj_semanal}, com duas diferenças de escala. O \textit{setup} é cobrado pelos tempos efetivos das matrizes de troca, $\tau^{\mathrm{F}}_{i'i}$ e $\tau^{\mathrm{C}}_{kk'}$, em vez do tempo médio $\tau_i$ que o nível semanal usa por não conhecer a sequência; e o carregamento de estoque é rateado pelos $|S|$ turnos da semana, não pelas 52 semanas do ano. A violação da meta aparece uma única vez, ao fim do horizonte, porque a meta que desce do nível semanal é um estado terminal e não uma exigência turno a turno, e é cobrada à mesma margem $m_i$ das demais faltas: entrar a semana seguinte com menos estoque do que o plano previu é falta adiada, não uma categoria de custo à parte.

As restrições \eqref{eq:forma_unica} a \eqref{eq:forma_da_posicao} definem a estrutura do turno, de fora para dentro: no máximo uma forma montada por máquina, a máquina só produz se tiver forma montada, toda posição de um turno ativo recebe exatamente um SKU, e esse SKU tem de pertencer à forma montada. É essa cadeia que reduz a sequência interna a uma sequência de cores dentro de um mesmo balão.

As trocas são detectadas por linearização de conjunção, com a mesma forma nos dois tipos: a variável de troca é forçada a 1 quando os dois estados que a definem valem 1 ao mesmo tempo, e o objetivo a empurra de volta a 0 no resto dos casos. A troca de forma \eqref{eq:troca_forma} compara a montagem em turnos consecutivos; a troca de cor \eqref{eq:troca_cor} compara posições consecutivas, e a versão para $n = 1$ compara a primeira posição do turno com a cor herdada, $w^{\mathrm{C}}_{kj,s-1}$. Um par ausente da matriz DE-PARA não gera variável de troca: em seu lugar entra a proibição \eqref{eq:transicao_proibida}, que impede que as duas cores se sucedam. Como uma troca de forma leva a máquina a um SKU de outro balão, nenhum par de mesma forma existe naquela fronteira e nenhum \textit{setup} de cor é cobrado junto, porque a troca de forma o subsume.

As equações \eqref{eq:estado_cor} a \eqref{eq:carryover_cor} implementam o \textit{carryover} de estado. A cor ao fim do turno é a da última posição ocupada, e existe exatamente uma cor montada por forma montada. Quando a máquina não produz no turno, $u_{js} = 0$ e as restrições \eqref{eq:carryover_forma} e \eqref{eq:carryover_cor}, escritas em valor absoluto por concisão e implementadas como o par de desigualdades correspondente, forçam forma e cor a permanecerem as do turno anterior. É o que impede que um turno ocioso apague a configuração da máquina e permita atravessar uma troca sem pagá-la, e é a contrapartida formal do mecanismo de estado descrito por \citeonline{soxgao1999}.

A capacidade \eqref{eq:capacidade_turno} tem a mesma forma da capacidade semanal \eqref{eq:capacidade_semanal}, com produção convertida em horas pela produtividade e as horas de \textit{setup} somadas ao lado esquerdo; o que muda é que agora há dois termos de \textit{setup} em vez de um. É uma restrição dura, e não um relatório de excedente: uma ordem de produção não atravessa a fronteira do turno. O lote mínimo \eqref{eq:lote_minimo_turno} repete a forma de \eqref{eq:lote_minimo}, com uma diferença: ele é imposto apenas quando a posição \emph{inicia} uma corrida de cor, porque o termo $y_{kjsn} - y_{kjs,n-1}$ vale 1 nesse caso e 0 quando a cor continua, o que evita exigir o mínimo de novo a cada posição de uma mesma corrida.

O bloco \eqref{eq:balanco_turno} a \eqref{eq:previsao_turno} espelha o balanço semanal \eqref{eq:balanco_semanal} no nível de SKU e turno, com uma diferença deliberada: as entregas são explícitas. Onde o nível semanal deduz a carteira entregue de um termo composto, aqui $e^{o}_{ks}$ e $e^{d}_{ks}$ nomeiam separadamente o que foi entregue contra pedido firme e contra previsão, e o que sobra alimenta o estoque. Essa separação não altera o ótimo, já que é implicada pelo balanço, mas torna a natureza de cada quilograma entregue uma saída do modelo, e não uma atribuição feita a posteriori.

\subsection{Coordenação entre os níveis}

O contrato entre os dois modelos é a restrição \eqref{eq:meta_final}: o estoque do balão $i$ ao fim do último turno deve alcançar a meta $\bar{I}_i$ que o nível semanal planejou para o fim daquela semana, e a violação $g_i$ é penalizada. É a única informação que desce, e ela é suave.

A escolha é deliberada e se opõe à alternativa de fixar as janelas de produção. Passar do nível semanal a alocação máquina-balão já resolvida garantiria consistência automática entre os planos, mas ao custo de tornar o nível operacional incapaz de reagir: qualquer prazo da carteira que não coubesse na janela herdada viraria atraso, sem possibilidade de realocar. Com a meta de estoque como único vínculo, o nível operacional decide livremente máquina, turno e sequência, e a consistência com o plano semanal é cobrada pelo estado terminal em vez de pelo caminho. É o esquema de metas com folga penalizada que \citeonline{bitranhaashax1981} formalizam para a desagregação em horizonte rolante, e a folga $g_i$ é o que impede que uma meta inatingível propague inviabilidade para o nível de baixo, no sentido apontado por \citeonline{schneeweiss2003}.

\subsection{Premissas}

As premissas a seguir delimitam o alcance das duas formulações.

\textbf{A demanda é determinística.} Carteira e previsão são dadas como conhecidas em cada execução, e a previsão vem do modelo da Seção~\ref{sec:modelo_demanda}, sem atualização dentro de uma execução. \textbf{A carteira é a parcela confirmada da demanda total}, de modo que a previsão entra líquida, isto é, previsão total menos carteira, e as duas correntes não se somam nem se sobrepõem no balanço.

\textbf{O \textit{setup} do nível semanal é independente de sequência.} Com \textit{bucket} de uma semana, a ordem de produção interna é decisão do nível operacional; $\tau_i$ é o tempo médio de entrada no balão $i$, calculado sobre a matriz de troca de formas.

\textbf{Uma máquina mantém uma forma por turno} e \textbf{ordens de produção não atravessam o turno}. Trocas de forma ocorrem apenas na fronteira de turno e consomem horas do turno que se inicia, enquanto produção e \textit{setups} de um turno cabem nas horas daquele turno. As duas premissas se sustentam pela escala dos tempos: uma troca de forma consome parcela relevante de um turno de oito horas.

\textbf{O estado da máquina sobrevive a turnos ociosos.} Forma e cor montadas permanecem enquanto a máquina não produz, e a troca seguinte é cobrada a partir desse estado. \textbf{O horizonte operacional é a primeira semana do plano}, e as demais semanas permanecem no nível semanal até a re-execução seguinte, em horizonte rolante.

\textbf{O \textit{setup} de forma é independente de máquina.} A matriz $\tau^{\mathrm{F}}_{i'i}$ é comum a todas as máquinas; a diferenciação entre equipamentos está no custo horário $c_j$. \textbf{A troca de forma subsume a troca de cor}, porque ao remontar a forma a máquina é preparada para a primeira cor do novo balão sem \textit{setup} adicional de cor.

\textbf{A taxa de carregamento é linear no tempo.} A taxa anual $\theta$ é rateada uniformemente pelas 52 semanas do ano e, no nível operacional, pelos turnos da semana. \textbf{A perecibilidade do látex não é modelada}: o recebimento diário e a capacidade dos tanques são tratados como não restritivos, e a decisão de produção é limitada por capacidade de moldagem, não por disponibilidade de composto.

\newpage

\bibliographystyle{abntex2-alf}
\bibliography{referencias}   

\end{document}
